\diarytitle{0073}{放物型2階線形偏微分方程式の標準形と初期値問題（$u_{xx} - 2u_{xy} + u_{yy} = 0$）}

方針: 判別式 $D=B^2-AC$ が $0$ となる放物型の場合、特性方程式
\[
\frac{dy}{dx} = \frac{B}{A}
\]
は重根であり特性曲線は1本しか得られない。その特性曲線を $\xi$ にとり、$\xi,\eta$ が独立になるように $\eta$ を任意に選ぶ（例えば $\eta=y$）と、標準形 $u_{\eta\eta}=0$ が得られる。これを積分して一般解 $u=f(\xi)+\eta\,g(\xi)$ を求めた後、初期条件から $f,g$ を定める。

$A=1,\ B=-1,\ C=1$ であるから
\[
D = B^{2}-AC = 1-1 = 0 \quad (\text{放物型})
\]
特性方程式は
\[
\frac{dy}{dx} = \frac{B}{A} = -1
\]
（重根なので特性曲線は1本）。積分して $y+x=c$。そこで
\[
\xi = y+x, \qquad \eta = y
\]
（$\eta$ は $\xi,\eta$ が独立になるように任意に選んだもの）とおく。逆に $x=\xi-\eta,\ y=\eta$。連鎖律より
\[
u_x = u_\xi, \qquad u_y = u_\xi + u_\eta
\]
\[
u_{xx}=u_{\xi\xi}, \qquad u_{xy}=u_{\xi\xi}+u_{\xi\eta}, \qquad u_{yy}=u_{\xi\xi}+2u_{\xi\eta}+u_{\eta\eta}
\]
元の方程式に代入すると
\[
u_{\xi\xi} - 2(u_{\xi\xi}+u_{\xi\eta}) + (u_{\xi\xi}+2u_{\xi\eta}+u_{\eta\eta})
= (1-2+1)u_{\xi\xi} + (-2+2)u_{\xi\eta} + u_{\eta\eta} = u_{\eta\eta}
\]
よって標準形は
\[
u_{\eta\eta} = 0
\]
これを $\eta$ について2回積分すると
\[
u = f(\xi) + \eta\,g(\xi)
\]
（$f,g$ は任意関数）。$\xi=y+x,\ \eta=y$ に戻すと一般解は
\[
u(x,y) = f(x+y) + y\,g(x+y)
\]
次に初期条件を課す。$u_y = f'(x+y) + g(x+y) + y\,g'(x+y)$ であるから
\[
u(x,0) = f(x) = x^{2}
\]
\[
u_y(x,0) = f'(x) + g(x) = 0 \quad\Longrightarrow\quad g(x) = -f'(x) = -2x
\]
これらを一般解に代入して、$f(x+y)=(x+y)^2$, $g(x+y)=-2(x+y)$ より
\[
u(x,y) = (x+y)^{2} - 2y(x+y)
\]
展開すると $(x+y)^2 - 2y(x+y) = x^2+2xy+y^2-2xy-2y^2 = x^2 - y^2$ なので
\[
\boxed{\; u(x,y) = x^{2} - y^{2} \;}
\]
（検算: $u=x^2-y^2$ より $u_{xx}=2,\ u_{xy}=0,\ u_{yy}=-2$ で $u_{xx}-2u_{xy}+u_{yy}=2-0-2=0$ を満たす。また $u(x,0)=x^2$、$u_y=-2y$ より $u_y(x,0)=0$ もともに満たす。）
